\section{Recurrence, ergodic and entropy}

  Recall that for two measure spaces $(X, \calF, \mu)$ and $(Y, \calG, \nu)$, a measurable map $f: X \to Y$ is called \emph{measure-preserving} if for every $A \in \calG$, we have $\mu(f^{-1}(A)) = \nu(A)$.

  \begin{example}
    The power map $S^1 \to S^1, z \mapsto z^n$ is measure-preserving with respect to the Lebesgue measure on the unit circle.
  \end{example}

  We fix a probability space $(X, \calF, \mu)$.


\subsection{Recurrence}

  \begin{theorem}[Poincar\'e Recurrence Theorem]
    Let $T: X \to X$ be a measure-preserving transformation. 
    For every measurable set $A$, we have 
    \[
    \mu\left(\{x \in A : T^n x \in A \text{ for infinitely many } n \in \mathbb{N}\}\right) = \mu(A).
    \]
    In other words, for almost every point in $A$, its orbit under $T$ returns to $A$ infinitely often.
    % \Yang{}
  \end{theorem}
  \begin{proof}
    \Yang{}
  \end{proof}

  \begin{example}
    \Yang{}
  \end{example}




\subsection{Ergodic}

  \begin{definition}
    A measurable transformation $T: X \to X$ is called \emph{ergodic} if for every $A \in \calF$ such that $T^{-1}(A) \aeeq A$, we have $\mu(A) = 0$ or $\mu(A) = 1$. 
  \end{definition}

  Recall that our $X$ is a probability space.
  In other words, the only $T$-invariant sets are trivial (i.e., empty set and the whole space up to measure zero).

  \begin{theorem}[Birkhoff's Ergodic Theorem]
    Let $T: X \to X$ be a measure-preserving transformation and $f \in L^1(X, \calF, \mu)$. 
    \begin{enumerate}
      \item Then for $\mu$-almost every $x \in X$, the limit
        \[
        \lim_{N \to \infty} \frac{1}{N} \sum_{n=0}^{N-1} f(T^n x)
        \]
        exists.
      \item Denote this limit by $\tilde{f}(x)$, then $\tilde{f}$ is integrable and $T$-invariant.
      \item If $T$ is ergodic, then $\tilde{f} = \int_X f \, d\mu$ is constant almost everywhere.
    \end{enumerate}
  \end{theorem}
  \begin{proof}
    
    \Yang{}
  \end{proof}


  Let $M(\calF, T)$ be the set of all $T$-invariant probability measures on $(X, \calF)$.
  Note that $M(\calF, T)$ is a convex set, that is, if $\mu_1, \mu_2 \in M(\calF, T)$ and $0 < t < 1$, then $t \mu_1 + (1-t) \mu_2 \in M(\calF, T)$.

  \begin{theorem}
    The ergodic measure in $M(\calF,T)$ are exactly the extreme points of $M(\calF,T)$.
    % \Yang{}
  \end{theorem}
  \begin{proof}
    \Yang{}
  \end{proof}



\subsection{Entropy}

  A \emph{partition} of $X$ is a collection $\calA = \{A_i\}_{i \in I}$ of disjoint measurable subsets of $X$ such that $\bigcup_{i \in I} A_i = X$. 
  For two partitions $\calA$ and $\calB$, the \emph{join} of $\calA$ and $\calB$ is defined as
  \[
  \calA \vee \calB = \{A \cap B : A \in \calA, B \in \calB\}.
  \]
  We say that $\calA$ is \emph{finer} than $\calB$ if for every $A \in \calA$, there exists $B \in \calB$ such that $A \subseteq B$, denoted by $\calA \leq \calB$.
  This is equivalent to saying that $\calA \vee \calB = \calA$.
  For every $x \in X$, we denote by $\calA(x)=A_x$ the unique element of $\calA$ containing $x$.


  \begin{definition}[Information function]
    Let $\calA = \{A_i\}_{i \in I}$ be a countable partition of $X$. The \emph{information function} of $\calA$ is defined as
    \[
    I_\calA(x) = -\log \mu(A_x),
    \]
    where $A_x$ is the unique element of $\calA$ containing $x$.
    % \Yang{}
  \end{definition}


  \begin{definition}[Entropy of a partition]
    Let $\calA = \{A_i\}_{i \in I}$ be a countable partition of $X$. The \emph{entropy} of $\calA$ is defined as
    \[
    H(\calA) = -\sum_{i \in I} \mu(A_i) \log \mu(A_i),
    \]
    where we adopt the convention that $0 \log 0 = 0$.
  \end{definition}

  The information function and the entropy of a partition are related by the following formula:
  \[ H(\calA) = \int_X I_\calA(x) \, d\mu(x). \]

  \begin{definition}[Conditional entropy]
    Let $\calA$ and $\calB$ be two countable partitions of $X$. 
    The \emph{conditional entropy} of $\calA$ given $\calB$ is defined as
    \[
      H(\calA | \calB) = \sum_{B \in \calB} \mu(B) H(\calA | B),
    \]
    where
    \[
      H(\calA | B) = -\sum_{A \in \calA} \frac{\mu(A \cap B)}{\mu(B)} \log \frac{\mu(A \cap B)}{\mu(B)}.
    \]
    %  \Yang{}
  \end{definition}

  We have 
  \[ H(\calA | \{X\}) = H(\calA). \]

  \begin{example}
    \Yang{}
  \end{example}

  Here we list some properties of the entropy and the conditional entropy.

  \begin{proposition}
    Let $\calA, \calB, \calC$ be three countable partitions of $X$. 
    \begin{enumerate}
      \item $H(\calA \vee \calB | \calC) = H(\calA | \calC) + H(\calB | \calA \vee \calC)$. 
        % In particular, $H(\calA \vee \calB) = H(\calA) + H(\calB | \calA)$.
      \item $ \calA \leq \calB \implies H(\calA | \calC) \leq H(\calB | \calC)$.
      \item $ \calB \leq \calC \implies H(\calA | \calB) \geq H(\calA | \calC)$.
      \item $H(\calA \vee \calB | \calC) \leq H(\calA | \calC) + H(\calB | \calC)$.
      \item $H(\calA | \calC) \leq H(\calA | \calB) + H(\calB | \calC)$.
    \end{enumerate}
  \end{proposition}
  \begin{proof}
    
    \Yang{May I omit them?}
  \end{proof}



  Let $T: X \to X$ be a measure-preserving transformation. 
  For a partition $\calA$ of $X$, we can define the partition $T^{-1} \calA$ as
  \[ T^{-1} \calA = \{T^{-1} A : A \in \calA\} \]
  and the partition $\calA^n$ as
  \[ \calA^n = \bigvee_{i=0}^{n-1} T^{-i} \calA. \]


  \begin{definition}[Entropy of an endomorphism]
    The \emph{entropy of $T: X \to X$ with respect to a partition $\calA$} is defined as
    \[ h_\mu(T, \calA) = \lim_{n \to \infty} \frac{1}{n} H(\calA^n). \]
    % where $\calA^n = \bigvee_{i=0}^{n-1} T^{-i} \calA$.
    The \emph{entropy of $T$} is defined as
    \[ h_\mu(T) = \sup_{\calA} h_\mu(T, \calA), \]
    where the supremum is taken over all finite measurable partitions $\calA$ of $X$.    
  \end{definition}

  \begin{example}
      \Yang{}
  \end{example}

  \begin{proposition}
    Let $T: X \to X$ be a measure-preserving transformation and $\calA, \calB$ be two countable measurable partitions of $X$ or finite entropy. 
    \begin{enumerate}
      \item $h_\mu(T, \calA) \leq H(\calA)$.
      \item $h_\mu(T, \calA \vee \calB) \leq h_\mu(T, \calA) + H(\calB)$.
      \item $ \calA \leq \calB \implies h_\mu(T, \calA) \leq h_\mu(T, \calB)$.
      \item $h_\mu(T, \calA \vee \calB) = h_\mu(T, \calA) + h_\mu(T, \calB | \calA)$.
      \item $h_\mu(T, \calA) \leq h_\mu(T, \calB) + h_\mu(T, \calA | \calB)$.
    \end{enumerate}
    % \Yang{}
  \end{proposition}
  \begin{proof}
    \Yang{}
  \end{proof}

  \begin{proposition}
    Let $T: X \to X$ be a measure-preserving transformation.
    Then 
    \[ h(T^k) = k h(T) \]
    for every positive integer $k$.
    Furthermore, if $T$ is invertible, then 
    \[ h(T^{-1}) = h(T). \]
      % \Yang{}
  \end{proposition}
  \begin{proof}
    \Yang{}
  \end{proof}


  \begin{theorem}[Shannon-McMillan-Breiman Theorem]
    Let $T: X \to X$ be an ergodic measure-preserving transformation and $\calA$ be a countable measurable partition of $X$ with finite entropy. 
    Then for $\mu$-almost every $x \in X$, we have
    \[
      \lim_{n \to \infty} \frac{1}{n} I_{\calA^n}(x) = h_\mu(T, \calA).
    \]
  \end{theorem}
  \begin{slogan}
    The ``local entropy'' $\lim_{n \to \infty} \frac{1}{n} I_{\calA^n}(x)$ equals the ``global entropy'' $h_\mu(T, \calA)$ for a.e. point $x$.
  \end{slogan}
  \begin{proof}
    \Yang{}
  \end{proof}


\subsection{On Lebesgue space}

  